% Paste-ready replacement for the outdated mean-field simulation, robustness,
% and parameter-table text. Values are read from paper_v2/configs/*.yaml.

\paragraph{Mean-field models and simulation.}
We studied six independently initialized mean-field models corresponding to the
familiar and novel $+NO$, $+O$, and $-NO$ transitions.  Each model contained one
PyC and one PV cell and was initialized directly from a YAML parameter file; no
population sampling or parameter averaging was performed during this study.
The PyC and PV instantaneous activations were rectified and clipped to $[0,1]$.
Thus the somatic activation threshold was zero.  Their discrete-time updates
used $\alpha_y=0.05$ and $\alpha_p=0.5$.

The fast adaptation state followed the PyC rate with
$\alpha_{a_f}=0.2\alpha_y=0.01$.  The slow activity accumulator used to modulate
FF plasticity followed the PyC rate with
$\alpha_{a_s}=0.025\alpha_y=0.00125$.  Consequently
$\tau_{a_s}>\tau_{a_f}>\tau_y$.  The apical gain onset was
$\theta_g=0.05$; the apical-drive threshold $\theta_d$, gain ceiling $g$, and
PyC baseline-current standard deviation $\sigma_y$ depended on the transition
(Table~\ref{tab:mean-field-dynamics}).

For image $i$, the NO input was the pair
$(\mathbf{x}_i,\mathbf{c}_i)$, where $\mathbf{x}_i$ was one-hot and
\begin{equation}
 \mathbf{c}_i = \operatorname{perm}_i([1,0.6,0.6])/2.2.
\end{equation}
The same unit-sum contextual vector reached the apical compartment and the FB
plasticity rule.  The O input replaced $\mathbf{x}_i$ with zero.  A trial had
400 steps ($100$ steps/s): three seconds without input followed by one second
of stimulus.  Displayed traces contained the final one second before stimulus,
the one-second stimulus, and two seconds of subsequent ITI.  Familiarization
comprised seven presentations of each familiar image in a reproducibly shuffled
order.  Naive and expert responses were measured over 50 trials per image.
Plasticity was active only during familiarization.

The three plastic updates were
\begin{align}
 \Delta\mathbf{w}_{FF}
 &= -\eta_{FF}\zeta a_s^4 y\mathbf{x}\odot\mathbf{w}_{FF},\\
 \Delta\mathbf{w}_{FB}
 &= \eta_{FB}\mathbf{c}\odot(\mathbf{1}-\mathbf{w}_{FB}),\\
 \Delta w_{LAT}
 &= \eta_{LAT}yp(1-w_{LAT}).
\end{align}
PV feedforward weights and the PyC-to-PV weight were fixed.  The rates were
$\eta_{FF}=0.5$, $\eta_{FB}=0.0003$, and
$\eta_{LAT}=0.05$ per discrete step, with no run-time rescaling.  The slow
accumulator normalization was $\zeta=1$ and its exponent was four.  Feedback
potentiation was undampened and no feedback-specificity matrix was applied.

For the main figures, each model was standardized using one shared mean and
sample standard deviation calculated over the one-second pre-stimulus windows
of all naive NO and O probes of all three images.  Familiar transition responses
were averaged over the two familiar images; novel responses used the novel
image.  Expert perturbations set contextual input to zero (FB silencing), removed
the PV-to-PyC term from the divisive denominator (PV silencing), or combined the
two manipulations.

\subsection{Robustness analysis}
The robustness analysis varied one continuous parameter at a time around each
of the six YAML initializations.  Every scalar parameter was multiplied by
\begin{equation}
 \{0.05,0.1,0.25,0.5,0.75,1,1.25,1.5,2,3,5\}.
\end{equation}
The leak rates $\alpha_y$ and $\alpha_p$, which are fractions per step, were
capped at one.  Each complete weight family ($\mathbf{w}_{FF}$,
$\mathbf{w}_{FB}$, $\mathbf{w}_{PV}$, $w_{LAT}$, and $w_{pvLAT}$) was scaled
coherently by the same factors, and the scaled weights were clipped to the
allowed interval $[0,1]$.  Every individual weight coordinate was instead set
to the absolute values
\begin{equation}
 \{0.01,0.02,0.05,0.1,0.2,0.3,0.5,0.75,1\}
\end{equation}
and to its YAML value, spanning the allowed range independently of its initial
size.  All unperturbed parameters remained at their YAML values.  Seeds and the
binary context-reception mask were not perturbed.  The design contained 1,965
perturbed points across the six models.

Robustness simulations used ten measurement trials by default.  Common random
numbers were used across perturbations, with separate fixed seeds for training
and measurement.  Responses were standardized using the full three-second
pre-stimulus period of the native-image naive probes; this is the pre-specified
robustness normalization and differs from the one-second plotting baseline used
for the main traces.  A perturbed transition was retained when it remained in
the same rotated sector (sector boundaries at the two diagonals) and its vector
magnitude was at least $25\%$ of the corresponding unperturbed magnitude.
Zero-variance baselines produced invalid, non-retained points.

\begin{table*}[h!]
\centering
\small
\begin{tabular}{lll}
\hline
Symbol & Value & Meaning \\
\hline
$\alpha_y$ & 0.05 & PyC leak rate \\
$\alpha_p$ & 0.5 & PV leak rate \\
$\alpha_{a_f}$ & 0.01 & Fast adaptation rate \\
$\alpha_{a_s}$ & 0.0005 & Slow FF-activity accumulator rate \\
% $\theta_y$ & 0 & PyC and PV rectification threshold \\
$\kappa$ & 10 & Divisive normalization scale \\
$\sigma_p$ & 0.15 & PV baseline-current s.d. \\
$\theta_g$ & 0.1 & Apical gain threshold \\
$k$ & 10 & Apical-gain sigmoid steepness \\
$\eta_{FF}$ & 0.6 & FF depression rate per step \\
$\eta_{FB}$ & 0.0003 & FB potentiation rate per step \\
$\eta_{LAT}$ & 0.009 & PV-to-PyC potentiation rate per step \\
$\zeta$ & 1 & Slow-accumulator scale \\
$q$ & 1 & Slow-accumulator exponent \\
$\mathbf{x}_i$ & $\text{perm}_i([1,0,0])$ & Feedforward input \\
$\mathbf{c}_i$ & $\text{perm}_i([1,.6,.6])/2.2$ & Contextual input \\
$T$ & 400 & Steps per trial (100 steps/s) \\
$N_{train}$ & 7 & Training trials per familiar image \\
$N_{test}$ & 50 & Naive and expert trials per image \\
\hline
\end{tabular}
\caption{Parameters shared by all six mean-field models.}
\label{tab:mean-field-shared}
\end{table*}

\begin{table*}[h!]
\centering
\scriptsize
\setlength{\tabcolsep}{3pt}
\resizebox{\textwidth}{!}{%
\begin{tabular}{llccccccr}
\hline
Source & Transition & $\mathbf{w}_{FF}$ & $\mathbf{w}_{FB}$ & $\mathbf{w}_{PV}$
& $w_{LAT}$ & $w_{pvLAT}$\\
\hline
Familiar & $+NO$ & $(.04,.004,.004)$ & $(.025,.025,.025)$ & $(.1,.1,.1)$ & .025 & 0\\
Familiar & $+O$  & $(.05,.05,.05)$ & $(.05,.05,.05)$ & $(.6,.6,.6)$ & .05 & 0\\
Familiar & $-NO$ & $(.06,.06,.06)$ & $(.006,.006,.006)$ & $(.1,.1,.1)$ & .025 & 0\\
Novel & $+NO$ & $(.02,.02,.02)$ & $(.025,.025,.025)$ & $(.1,.1,.1)$ & .025 & 0\\
Novel & $+O$  & $(.05,.05,.05)$ & $(.05,.05,.05)$ & $(.6,.6,.6)$ & .05 & 0\\
Novel & $-NO$ & $(.05,.05,.05)$ & $(.005,.005,.005)$ & $(.5,.5,.5)$ & .05 & 0\\
\hline
\end{tabular}}
\caption{Initial synaptic parameters of the six mean-field models.}
\label{tab:mean-field-weights}
\end{table*}

\begin{table}[h!]
\centering
\small
\begin{tabular}{llccc}
\hline
Source & Transition & $\theta_d$ & $g$ & $\sigma_y$ \\
\hline
Familiar & $+NO$ & .2 & 18 & .15 \\
Familiar & $+O$  & .14 & 8 & .15 \\
Familiar & $-NO$ & .2 & 5 & .15 \\
Novel & $+NO$ & .14 & 12 & .15 \\
Novel & $+O$  & .14 & 8 & .15 \\
Novel & $-NO$ & .2 & 5 & .15 \\
\hline
\end{tabular}
\caption{Transition-specific neural-dynamics parameters. All unlisted values
are shared as in Table~\ref{tab:mean-field-shared}.}
\label{tab:mean-field-dynamics}
\end{table}
